\documentclass{article}
\usepackage[T1]{fontenc}
\usepackage{lmodern}
\usepackage{graphicx}
\usepackage{amsmath,amssymb}
\usepackage[a4paper,margin=2cm]{geometry}
\usepackage[absolute,overlay]{textpos}
\setlength{\TPHorizModule}{1mm}
\setlength{\TPVertModule}{1mm}
\usepackage[indonesian]{babel}
\usepackage{hyperref}
\setlength{\emergencystretch}{2em}
% unit: Halaman 14

\begin{document}

% === Halaman 14 (python/native) ===

• Contoh di dalam Vigenere Cipher, cipherteks setiap huruf plainteks dihitung 

dengan:

               𝐶= (𝑃+ 𝐾) 𝑚𝑜𝑑 26

• Penyerang memilih plaintext P = AAAAAAA… (= 000000…).

• Maka cipherteks yang keluar, C, pasti sama dengan K    (C = K).

• Dengan mengetahui K, cipherteks lain bisa langsung dicoba didekripsi dengan K 

tersebut.

14

\end{document}
